\section{Introduction}
\label{sec:introduction}

The abundance conjecture predicts that a nef canonical divisor on a
projective variety with mild singularities is semiample: some positive
Cartier multiple is generated by its global sections.  Such sections
turn numerical information about the canonical divisor into a morphism
that records the canonical geometry of the variety.  For a nef divisor
$L$ on a threefold, numerical dimension one means
$L\cdot H^2>0$ and $L^2\cdot H=0$ for an ample Cartier divisor $H$.
We prove this case of abundance for terminal threefolds in
characteristic greater than three.

\begin{theorem}\label{thm:main}
Let $X$ be a normal projective $\Q$-factorial threefold over an
algebraically closed field $k$ of characteristic $p>3$.  Suppose that
$X$ has terminal singularities and that $K_X$ is nef with
$\nu(K_X)=1$.  Then some positive Cartier multiple of $K_X$ has at
least two independent global sections.  In particular,
$\kappa(X,K_X)=1$ and $K_X$ is semiample.
\end{theorem}

No assumption is made on the first Betti number, irregularity,
or Albanese variety.  Thus Theorem~\ref{thm:main} resolves the
numerical-dimension-one case of the abundance conjecture under the
stated terminal and characteristic hypotheses.

Geometrically, the theorem gives a morphism $f:X\to C$ with connected
fibers to a normal projective curve and an ample Cartier divisor $A$
on $C$ such that $mK_X\sim f^*A$ for some $m>0$.  Indeed, the image
of a morphism defined by a globally generated pluricanonical multiple
is a curve, since $\nu(K_X)=1$.  Its Stein factorization gives $f$;
the finite pullback to $C$ of a hyperplane divisor on the image gives $A$.

\subsection{Context}

Over the complex numbers, Miyaoka proved the numerical-dimension-one
case of threefold abundance by studying effective pluricanonical
divisors, primitive connected cycles, and their infinitesimal
neighborhoods \cite{miyaoka1988}.  Kawamata's abundance theorem for
minimal threefolds includes another proof of this case using the log
minimal model program and formal thickenings
\cite[Section~4]{kawamata1992}.  The present argument develops this
canonical-cycle framework in positive characteristic.
More recently, Liu--Xu developed higher-dimensional
numerical-dimension-one abundance over $\mathbb C$ under nonvanishing
hypotheses \cite{liuxu2025}.  The ground field in the present theorem
requires the positive-characteristic inputs described next.

The positive-characteristic minimal model program supplies the
birational operations needed to study abundance for threefolds.
Over algebraically closed fields of characteristic greater than five,
Hacon--Xu established standard-coefficient flips and minimal models
for projective $\Q$-factorial terminal threefolds with pseudoeffective
canonical divisor \cite[Theorems~1.1 and~1.2]{haconxu2015}.
Birkar proved flips for $\Q$-factorial dlt threefold pairs and log
minimal models for projective klt threefold pairs with pseudoeffective
log canonical divisor \cite[Theorems~1.1 and~1.2]{birkar2016}.
The relative log minimal model program in characteristic greater than
five was established by Hashizume, Nakamura and Tanaka
\cite{hnt2020}; Hacon and Witaszek treated the characteristic-five
case \cite{hw2022}.

In characteristic greater than five,
Xu--Zhang proved nonvanishing for terminal minimal threefolds
\cite{xuzhang2019}; Witaszek obtained nonvanishing for klt threefold
pairs and abundance in nef dimension at most two through a canonical
bundle formula \cite{witaszek2021}.  Zhang proved abundance for
minimal threefolds with nontrivial Albanese map \cite{zhang2020}.
Xu proved nonvanishing for pseudoeffective log
canonical threefold pairs in characteristic greater than three and
semiampleness when the nef log canonical divisor has positive Kodaira
dimension \cite[Theorems~1.4 and~1.6]{xu2024}.  His theorem for nef
dimension at most two \cite[Theorem~1.8]{xu2024} and his subsequent
numerical-dimension-two abundance theorem \cite[Theorem~1.1]{xu2026}
leave a specific issue in the present setting: an effective canonical
divisor of numerical dimension one could a priori still have Kodaira
dimension zero and maximal nef dimension.  Theorem~\ref{thm:main}
excludes that possibility.  Here nef dimension is the dimension of the
nef reduction, and maximal nef dimension means that no curve of
canonical degree zero passes through a very general point.  It is a
different invariant from numerical dimension.  The new step is the
production of a second pluricanonical section; the existence of a first
section and the passage from positive Kodaira dimension to semiampleness
are Xu's results.

Our tangent-sheaf argument follows the positive-characteristic
foliation method appearing in Xu--Zhang \cite[Section~2]{xuzhang2019}.
It uses the Bogomolov inequality for strongly semistable sheaves with
respect to a nef tuple, in Langer's formulation
\cite[Theorem~3.2]{langer2004}; the corrections in
\cite{langeradd2004} preserve that theorem.  Surface abundance enters
through abundance for projective semi-log-canonical surfaces, proved
by Tanaka over algebraically closed fields \cite[Theorem~0.1]{tanaka2016}
and by Posva over arbitrary fields \cite[Theorem~3.1]{posva2024}.

The descent mechanism is classical Cartier theory.  In particular,
Katz's proof of Cartier descent supplies the horizontal projector and
truncated Taylor expansion \cite[Theorem~5.1]{katz1970}.
We establish the completed-ring and bounded-pole statements that allow
this algebra to be used on the punctured formal neighborhood of a
nonreduced cycle.  The three-row jet calculation and its uniform pole
control are the bridge between the canonical-cycle geometry and this
descent argument.

\subsection{The argument}

After nonvanishing and the birational preparation in
Section~\ref{sec:preparation}, the hypothetical
maximal-nef-dimension case produces a smooth projective threefold $V$
and a connected effective nef Cartier cycle $D$ with primitive
multiplicities.  The cycle is numerically trivial on each of its
components, its normal line $\cO_D(D)$ is torsion, and, on a neighborhood
of $D$, a $p$-power of $\omega_V$ has a divisor supported on $D$.
The preparation proves
torsion on the full nonreduced cycle, which is needed in the subsequent
calculation.

Section~\ref{sec:slopes} then proves strong semistability of the tangent
and cotangent bundles for the nef degree $c_1(-)\cdot H\cdot D$.
The foliation quotient of a destabilizing sheaf would produce
$D$-trivial curves through very general points.  The resulting
semistability gives $c_2(V)\cdot D\geq0$ and uniform bounds on polar
parts, which are the inputs to the jet calculation.

The central calculation in Section~\ref{sec:jets} concerns finite jets
along $D$, that is, sections on successive Cartier thickenings.  It keeps all
three cohomological rows of the pole filtration and compares their
ranks over the periods $rp^e$, where $r$ is the order of
$\cO_D(D)$.  The scalar rows identify the completed canonical line with
the line of an integral divisor supported on $D$, and force both
$c_2(V)\cdot D=0$ and the
precise limiting ratio $p/(p+1)$ for the last trivialization widths.
These facts control the positions and lengths of differentials for
vector bundles as well.

Section~\ref{sec:laurent} turns the calculation into a
formal-neighborhood criterion,
Theorem~\ref{thm:laurent-vanishing}: a bundle satisfying an Euler
identity, bounded polar parts, and a concentration condition for
leading jets has vanishing first Laurent cohomology.  A bound on the
source cuts of incoming differentials ensures that the constructed
primitives actually converge with bounded poles.  We prove this
criterion directly for finite covers and completed lattices.
The criterion applies to one-forms and exact forms.

Section~\ref{sec:cartier} uses these vanishings to make every line
bundle on the punctured formal neighborhood a $p$th power, by Cartier
descent.  On the other hand, intersection with $H$ on the components
of $D$ defines a nonzero integer-valued homomorphism on its Picard
group.  Such a homomorphism cannot exist on a $p$-divisible group.
The contradiction proves Theorem~\ref{thm:main}.

Figure~\ref{fig:proof-map} records the two conclusions about the same
Picard group that give the contradiction.  In the figure,
$\cR=\widehat{\cO}_V(*D)$ denotes functions on the completion along
$D$ with a finite pole bound; its Picard group consists of locally free
rank-one $\cR$-modules on $|D|$.
\begin{figure}[ht]
\centering
\small
\begin{tikzcd}[column sep=small,row sep=large]
& \text{Prepared cycle }(V,D)
  \arrow[dl,"\text{Sections~\ref{sec:slopes}--\ref{sec:laurent}}"']
  \arrow[dr,"\text{Section~\ref{sec:cartier}}"] & \\
\begin{gathered}\text{Slope bounds and finite jets}\\
            \text{Laurent cocycles have primitives}\end{gathered}
  \arrow[d,"\text{Cartier descent}"']
&& \begin{gathered}\text{Orders of Laurent units}\\
             \text{Intersection with }H\end{gathered}
  \arrow[d] \\
\begin{gathered}\Pic(|D|,\cR)\\\text{is }p\text{-divisible}\end{gathered}
  \arrow[dr]
&& \begin{gathered}\overline\Phi:\Pic(|D|,\cR)\to\Z\\
             \overline\Phi(\cO(H))=DH^2>0\end{gathered}
  \arrow[dl] \\
& \text{Contradiction} &
\end{tikzcd}
\caption{The maximal-nef-dimension assumption produces a cycle for
which Laurent cohomology forces Picard divisibility, while intersection
theory supplies a nonzero integral degree.}
\label{fig:proof-map}
\end{figure}

All arguments take place in characteristic $p$.  The proof uses no
lifting to characteristic zero and no orbifold Iitaka theorem.

\subsection{Conventions}

Varieties are integral and all divisors are rational Cartier divisors
when a pullback or intersection requires this, unless they are
explicitly declared integral Cartier divisors.  On a smooth variety
we identify Cartier and Weil divisors.  Linear and numerical
equivalence are denoted by $\sim$ and $\equiv$, with a subscript
$\Q$ for rational linear equivalence.  For a nef divisor $L$, its
nef dimension $n(L)$ is the dimension of its nef reduction; in
particular, $n(L)=\dim X$ means that no $L$-trivial curve passes
through a very general point.  We use an uncountable algebraically
closed field when making this formulation, and descend the final
section-space statement to the original field.

If $E$ is a vector bundle on a smooth threefold $V$, write
$E^!=E^\vee\otimes\omega_V$.  All completions below are along an
effective Cartier divisor, and all Laurent sheaves are considered on
its Zariski support.  The term \emph{width} refers to the number of
successive Cartier layers in a jet.
